\documentclass[11pt]{article}

\topmargin=-0.5in
\oddsidemargin=0in
\evensidemargin=0in
\textwidth=6.5in
\textheight=9.0in

\usepackage{setspace}
\usepackage{graphicx}

\usepackage{amsmath}
\usepackage{amssymb}
\usepackage{amsthm}
\usepackage{url}

\usepackage{setspace}

\input{Qcircuit}

% \noindent {\sc Slide}: Jones-Gomard-Sestoft book cover.

% \mbox{} \\ \hspace*{1cm}
%    [Our experiments confirm our theory] \\

\title{Quantum Algorithms \\
       Clifford+T is Universal}
\author{Jens Palsberg}
\date{June, 2026}

% Title: Towards a universal quantum programming language.
% Goal: 2,500 words.
% Almost every quantum computer right now has its own programming language. 
% How may this change for the future quantum computers and what are 
% the unique challenges in making a universal quantum computing language?
%
% Role model: http://xrds.acm.org/current-issue.cfm

\def\({\begin{eqnarray*}}
\def\){\end{eqnarray*}}

\newtheorem{theorem}{Theorem}
\newtheorem{corollary}[theorem]{Corollary}
\newtheorem{lemma}[theorem]{Lemma}

\def\phi{\varphi}
\def\imp{\rightarrow}
\def\implies{\Rightarrow}

\def\ket#1{|#1\rangle}
\def\bra#1{\langle#1|}
\def\inner#1#2{\langle{#1}|{#2}\rangle}
\def\outer#1#2{\ket{#1}\bra{#2}}
\def\eigenvalues#1{\mbox{\sf Eigenvalues}(#1)}

\def\diag#1{\mbox{\sf Diag}{(#1)}}
\def\diagtwo#1#2{\diag{#1,#2}}

\def\C#1{\text{\sf C}{(#1)}}
\def\CC#1{\text{\sf CC}{(#1)}}


\def\iff{\Longleftrightarrow}

\begin{document}
\maketitle

% Example
% Story
% Data
% Quote
% Peer Testimony

\newpage

\rightline{Quantum Algorithms, by Jens Palsberg}
\rightline{Clifford+T is Universal; 100 minutes; June, 2026} 

\vspace*{-1.2cm}

\section*{Outline}

\paragraph{Hook:} 
A quantum computer implements a universal gate set.
Intuitively, this means that the computer can implement any unitary. 
But technically, it means that the computer can approximate any unitary.
What is behind this idea and how can we prove that 
a gate set can approximate any unitary?

\paragraph{Purpose:} Persuade you that Clifford+T is universal.

\paragraph{Preview:} \mbox{} \\
\hspace*{1cm}
1. We will state the theorem and devise a proof strategy.
\\
\hspace*{1cm}
2. We will first prove a version of the theorem that works with $R_z$ gates.
\\
\hspace*{1cm}
3. We will prove the full theorem by approximating the $R_z$ gates.

\paragraph{Transition to Body:} 
Let us begin with a look at what we are trying to prove.

\paragraph{Main Point 1:}
We will state the theorem and devise a proof strategy. \\
\hspace*{1cm}
   [We state the theorem] \\ 
\hspace*{1cm}
   [We will first prove a version of the theorem that uses an easy gate set] \\
\hspace*{1cm}
   [We state some known lemmas]

\paragraph{Transition to MP2:}
Now we move on to a proof for an easy gate set.

\paragraph{Main Point 2:}
We will first prove a version of the theorem that works with $R_z$ gates. \\
\hspace*{1cm}
   [We prove the simple version of the theorem] \\
\hspace*{1cm}
   [We need exponentially many gates to implement an $n$-qubit unitary in the worst case] \\
\hspace*{1cm}
   [We prove one of the known lemmas] 

\paragraph{Transition to MP3:}
Now we can finish the proof by focusing on approximation.

\paragraph{Main Point 3:}
We will prove the full theorem by approximating the $R_z$ gates. \\
\hspace*{1cm}
   [We list properties of approximation of unitaries] \\
\hspace*{1cm}
   [We show how to approximate $R_z$ gates] \\
\hspace*{1cm}
   [We prove the full version of the theorem] 

\paragraph{Transition to Close:} 
So there you have it.

\paragraph{Review:}
The key property of a universal gate set is that 
it can approximate any unitary.
Our proof strategy is to first prove the theorem for an easy gate set and then
extend to $\{\C{X}, H, T\}$.

\paragraph{Strong finish:}
Quantum computing has probabilities built in and
the idea of a universal gate set has approximation built in.
Quantum algorithms have no requirement to be exact; 
rather, they can try to get close.

\paragraph{Call to action:}
Implement the procedure for compiling any unitary into the easy gate set.

\newpage

\section*{Detailed presentation}

\paragraph{Hook:} 
A quantum computer implements a universal gate set.
Intuitively, this means that the computer can implement any unitary. 
But technically, it means that the computer can approximate any unitary.
What is behind this idea and how can we prove that 
a gate set can approximate any unitary?

\paragraph{Purpose:} Persuade you that Clifford+T is universal.

\paragraph{Preview:} \mbox{} \\
\hspace*{1cm}
1. We will state the theorem and devise a proof strategy.
\\
\hspace*{1cm}
2. We will first prove a version of the theorem that works with $R_z$ gates.
\\
\hspace*{1cm}
3. We will prove the full theorem by approximating the $R_z$ gates.

\paragraph{Transition to Body:} 
Let us begin with a look at what we are trying to prove.

\paragraph{Main Point 1:}
We will state the theorem and devise a proof strategy. \\
\mbox{} \\ \hspace*{1cm}
   [We state the theorem] \\ 
We say that two unitaries $U$ and $V$ are equal up to a global phase,
written $U \sim V$, if there exists a real number $\alpha$ such that
$U = e^{i\alpha} V$, where $e^{i\alpha}$ is the global phase.
%
We define the following projective Hilbert–Schmidt distance $d$ for
$n$-qubit unitaries $U, V$:
\(
d(U,V) &=& 
\sqrt{ 1 - \frac{ | Tr( U^\dagger V ) |^2 }{ N^2 } }
\) 
where $Tr(U)$ denotes the trace of $U$, and $N = 2^n$.
This distance is invariant under multiplication of either unitary by 
a global phase.
%
Let us define three rotation matrices:
\[
\begin{array}{rclclcl}
R_x(\theta) &=& 
e^{-i\theta X/2}
&=&
\cos(\theta/2)\ I - i\sin(\theta/2)\ X
&=&
\left(
\begin{array}{cc}
\cos(\theta/2)    & -i \sin(\theta/2) \\
-i \sin(\theta/2) & \cos(\theta/2)
\end{array}
\right)
\\
R_y(\theta) &=&
e^{-i\theta Y/2}
&=&
\cos(\theta/2)\ I - i\sin(\theta/2)\ Y
&=&
\left(
\begin{array}{cc}
\cos(\theta/2)    & -\sin(\theta/2) \\
\sin(\theta/2) & \cos(\theta/2)
\end{array}
\right)
\\
R_z(\theta) &=&
e^{-i\theta Z/2}
&=&
\cos(\theta/2)\ I - i\sin(\theta/2)\ Z
&=&
\left(
\begin{array}{cc}
e^{-i\theta/2} & 0 \\
0              & e^{i\theta/2}
\end{array}
\right)
\end{array}
\]
%
We use the term circuit to mean a product of unitaries,
possibly from a particular gate set.  
The gates $\{\C{X}, H, S\}$ are called Clifford gates,
while $T$ and $R_z(\theta)$ are not Clifford gates.
We say that a gate set is universal if,
for any unitary $U$ and any $\epsilon > 0$, 
we can use the gate set to define a circuit $C$
such that $d(U,C) < \epsilon$.
Theorem~\ref{thm:clifford-plus-t-is-universal} below states that 
$\{\C{X}, H, T\}$ is universal.

\begin{theorem}[\bf Clifford+T is universal]
\label{thm:clifford-plus-t-is-universal}
For any unitary $U$ and any $\epsilon > 0$, 
there exists a circuit $C$ built from $\{\C{X}, H, T\}$,
such that $d(U,C) < \epsilon$.
\end{theorem}

% \mbox{} \\ 
\hspace*{1cm}
   [We will first prove a version of the theorem that uses an easy gate set] \\
Let us state an easier version that has no need for approximation.

\begin{lemma}[\bf Clifford+R$_z$ is universal]
For any $n$-qubit unitary $U$,
there exists a circuit $C$ with 
$14 \times 4^{n-1} - 9 \times 2^{n-1}$ gates drawn from
$\{\C{X}, H, S, R_z(\theta)\}$ such that $U \sim C$.
\label{lem:clifford-plus-rx-is-universal}
\end{lemma}

Our proof strategy is to first prove 
Lemma~\ref{lem:clifford-plus-rx-is-universal}
and then approximate the $R_z(\theta)$ gates.

\newpage

\mbox{} \\ \hspace*{1cm}
   [We state some known lemmas] \\
Lemma~\ref{lem:clifford-plus-rx-is-universal-for-1-qubit-gates} below
% and
% Lemma~\ref{lem:clifford-plus-rx-is-universal-for-2-qubit-gates} 
% are special cases 
is a special case of Lemma~\ref{lem:clifford-plus-rx-is-universal},
for 1-qubit gates.
% for 1-qubit gates and 2-qubit gates.
Our strategy for proving 
Lemma~\ref{lem:clifford-plus-rx-is-universal}
% for an $n$-qubit unitary, 
% where $n \geq 2$,
is to do induction with
Lemma~\ref{lem:clifford-plus-rx-is-universal-for-1-qubit-gates} 
as the base case.

\begin{lemma}[\bf Nielsen and Chuang 2000]
For any 1-qubit unitary $U$,
there exists a circuit $C$ with
2 $H$ gates and 3 $R_z(\theta)$ gates
such that $U \sim C$.
\label{lem:clifford-plus-rx-is-universal-for-1-qubit-gates}
\end{lemma}

\begin{proof}
We can find 
real numbers $\alpha,\beta,\gamma,\delta$ and write:
$U =
e^{i\alpha}\,
R_z(\beta)\,
H\,
R_z(\gamma)\,
H\,
R_z(\delta)$.
\end{proof}

% \begin{lemma}[\bf Chang, Vala, Sastry, and Whaley 2003]
% We can implement any 2-qubit unitary using 
% 6 $\C{X}$ gates, 16 $H$ gates, 8 $S$ gates, and 15 $R_z(\theta)$ gates,
% up to a global phase.
% \label{lem:clifford-plus-rx-is-universal-for-2-qubit-gates}
% \end{lemma}
% 
% \begin{proof}
% Let $U$ be a 2-qubit unitary.
% We can find 1-qubit unitaries $V_1, V_2, V_3, V_4$ and
% real numbers $\theta_X, \theta_Y, \theta_Z$ and write:
% \(
% U
% &=&
% (V_1 \otimes V_2)\
% e^{i \theta_X (X \otimes X)}\
% e^{i \theta_Y (Y \otimes Y)}\
% e^{i \theta_Z (Z \otimes Z)}\
% (V_3 \otimes V_4),
% \\[1ex]
% e^{i \theta_X (X \otimes X)}
% &=&
% (H \otimes H)\
% \C{X}\
% (I \otimes R_z(2\theta_X))\
% \C{X}\
% (H \otimes H),
% \\[1ex]
% e^{i \theta_Y (Y \otimes Y)}
% &=&
% (SH \otimes SH)\
% \C{X}\
% (I \otimes R_z(2\theta_Y))\
% \C{X}\
% (HS^\dagger \otimes HS^\dagger),
% \\[1ex]
% e^{i \theta_Z (Z \otimes Z)}
% &=&
% \C{X}\
% (I \otimes R_z(2\theta_Z))\
% \C{X}.
% \)
% We use Lemma~\ref{lem:clifford-plus-rx-is-universal-for-1-qubit-gates} to 
% get counts for the four 1-qubit unitaries, and we write $S^\dag = SSS$.
% \end{proof}

Our next lemma states the Cosine-Sine decomposition,
also known as the CS-decomposition.

\begin{lemma} {\bf (Paige and Wei 1994)}
For one n-qubit gate $U$, there exist $2^{n-1}$ 1-qubit gates $R_y(\theta_0)$, $R_y(\theta_1)$, $\cdots$, and $R_y(\theta_{2^{n-1}-1})$, four $(n-1)$-qubit gates $P_0$, $P_1$, $Q_0$, and $Q_1$, and three n-qubit gates $P = \outer{0}{0} \otimes P_{0} + \outer{1}{1} \otimes P_{1}$, $R = \sum^{2^{n-1}-1}_{i=0}R_y(\theta_i) \otimes \outer{i}{i}$, and $Q= \outer{0}{0} \otimes Q_{0} + \outer{1}{1} \otimes Q_{1}$, 
such that
\[ 
\Qcircuit @C=0.5em @R=0.4em {
  & \qw & \multigate{2}{U} & \qw & &&  & & \qw & \gate{} \qwx[2]& \gate{R}      &\gate{} \qwx[2]&\qw\\
 &  &                &     & &=& &   &    &          &                 &               &   \\
 & \qw / & \ghost{U}        & \qw &  && & & \qw / & \gate{Q}       & \gate{} \qwx[-2]&\gate{P}       &\qw\\
}
\]
\label{cosinesine}
\end{lemma}

% \newpage

Our next lemma says that we can implement an $R_y(\alpha)$ gate via use of 
an $R_z(-\alpha)$ gate.

\begin{lemma} {\bf (Shende, Markov, and Bullock 2004)}
The following circuit equality is correct: 
$R_y(\alpha) = S^\dagger\ H\ R_z(-\alpha)\ H\ S$, which can be written
also in circuit notation:
\[ 
\Qcircuit @C=0.5em @R=0.4em {
 & \gate{R_y(\alpha)} & \qw & & = & &  & \gate{S} & \gate{H} & \gate{R_z(-\alpha)}      &\gate{H}    & \gate{S^\dag} & \qw 
}
\]
\label{ryrz}
\end{lemma}

\vspace*{-0.8cm}

\begin{proof}
We calculate:
%
\(
S^\dagger\ H\ R_z(-\alpha)\ H\ S
&=&
S^\dagger\ H\ [\cos(-\alpha/2)\ I - i\sin(-\alpha/2)\ Z]\ H\ S
\\
&=&
S^\dagger\ [\cos(-\alpha/2)\ I - i\sin(-\alpha/2)\ X]\ S
\\
&=&
[\cos(-\alpha/2)\ I - i\sin(-\alpha/2)\ (-Y)]
\\
&=&
\cos(\alpha/2)\ I - i\sin(\alpha/2)\ Y
\\
&=&
R_y(\alpha)
\)
In the second step, we use $H H = I$ and $H Z H = X$.
In the third step, we use $S^\dagger S = I$ and $S^\dagger X S = - Y$.
\end{proof}

\begin{lemma} {\bf (Shende, Bullock, and Markov 2005)}
For two $(n-1)$-qubit gates $U_{0}$ and $U_{1}$ and one n-qubit gate $U = \outer{0}{0} \otimes U_{0} + \outer{1}{1} \otimes U_{1}$, there exist $2^{n-1}$ 1-qubit gates $R_z(\alpha_0)$, $R_z(\alpha_1)$, $\cdots$, and $R_z(\alpha_{2^{n-1}-1})$, two $(n-1)$-qubit gates $P$ and $Q$, and one n-qubit gate $R = \sum^{2^{n-1}-1}_{i=0}R_{z}(\alpha_i) \otimes \outer{i}{i}$, such that
\[ 
\Qcircuit @C=0.5em @R=0.4em {
& \qw & \gate{} \qwx[2] & \qw &  && &  & \qw & \qw       & \gate{R}      &\qw            &\qw\\  
& & &     & & = & &  &             &        &         &               &   \\
& \qw / & \gate{U}       & \qw &  && &  & \qw / & \gate{Q}       & \gate{} \qwx[-2]&\gate{P}       &\qw\\
}
\]
\label{demultiplexing}
\end{lemma}

Now we can count the number of gates needed to implement 
a multiplexed $R_z(\alpha)$ gate.

\begin{lemma}
Let $R$ be the $n$-qubit gate 
$\sum_{i=0}^{2^{n-1}-1} R_z(\alpha_i) \otimes \outer{i}{i}$.
We can implement $R$ using 
$2^n - 2$ $\C{X}$ gates and 
$2^{n-1}$ $R_z(\theta)$ gates.
\label{rzcount}
\end{lemma}

We defer the proof of Lemma~\ref{rzcount} until later.

\paragraph{Transition to MP2:}
Now we move on to a proof for an easy gate set.

\newpage

\paragraph{Main Point 2:}
We will first prove a version of the theorem that works with $R_z$ gates. 

\mbox{} \\ \hspace*{1cm}
   [We prove the simple version of the theorem] \\
We prove Lemma~\ref{lem:clifford-plus-rx-is-universal},
which states that we can implement any unitary $U$ using Clifford+R$_z$ gates.

\begin{proof}[Proof of Lemma~\ref{lem:clifford-plus-rx-is-universal}]
Suppose $U$ is an $n$-qubit unitary.
We proceed by induction on $n$.

In the base case of $n=1$, we use 
Lemma~\ref{lem:clifford-plus-rx-is-universal-for-1-qubit-gates}.

In the induction step, suppose Lemma~\ref{lem:clifford-plus-rx-is-universal} 
holds for $n$, where $n\geq 2$.  
Consider the case of where $U$ is an $(n+1)$-qubit unitary.
We calculate as follows.
%
\begin{eqnarray}
% && 
\Qcircuit @C=0.5em @R=1em {
& \qw & \multigate{2}{U}  & \qw \\
& \qw & \ghost{U} & \qw \\
& / \qw & \ghost{U}  & \qw 
}\nonumber
% \\[0.4cm] &&
& \;\;\;\;\;\;\; &
\Qcircuit @C=0.5em @R=1em {
& \qw & \gate{} \qwx[1] & \gate{R}  & \gate{} \qwx[1] & \qw \\
\lstick{= \;\;\;\; } & \qw & \multigate{1}{Q} &  \gate{} \qwx[-1] & \multigate{1}{P} & \qw \\
& / \qw & \ghost{Q}  &  \gate{} \qwx[-1] & \ghost{P}  & \qw 
}\nonumber
% \\[0.4cm] &&
\;\;\;\;\;\;\; \;\;\;\;\;\;\;\,
\Qcircuit @C=0.5em @R=1em {
& \qw & \gate{} \qwx[1] & \gate{S} & \gate{H} & \gate{R_1} & \gate{H} &\gate{S^\dag} & \gate{} \qwx[1] & \qw \\
\lstick{= \;\;\;\; } & \qw & \multigate{1}{Q} & \qw & \qw &  \gate{} \qwx[-1] & \qw & \qw & \multigate{1}{P} & \qw \\
& / \qw & \ghost{Q} & \qw & \qw &  \gate{} \qwx[-1] & \qw & \qw & \ghost{P}  & \qw 
}\nonumber
% \\[0.4cm] &&
% \Qcircuit @C=0.5em @R=1em {
% & \qw & \gate{} \qwx[1] & \gate{S} & \gate{H} & \gate{R_1} & \gate{H} &\gate{S^\dag} & \qw & \gate{R_0} \qwx[1] & \qw & \qw \\
% \lstick{= \;\;\;\; } & \qw & \multigate{1}{Q} & \qw & \qw &  \gate{} \qwx[-1] & \qw & \qw & \multigate{1}{W} & \gate{} \qwx[-1] & \multigate{1}{V} & \qw \\
% & / \qw & \ghost{Q} & \qw & \qw &  \gate{} \qwx[-1] & \qw & \qw & \ghost{W} & \gate{} \qwx[-1] & \ghost{V}  & \qw 
% }\nonumber
\\[0.4cm] &&
\Qcircuit @C=0.5em @R=1em {
& \qw & \qw & \gate{R_2} \qwx[1] & \qw & \gate{S} & \gate{H} & \gate{R_1} & \gate{H} &\gate{S^\dag} & \qw & \gate{R_0} \qwx[1] & \qw & \qw \\
\lstick{= \;\;\;\; } & \qw & \multigate{1}{N} & \gate{} \qwx[-1] & \multigate{1}{M} & \qw & \qw &  \gate{} \qwx[-1] & \qw & \qw & \multigate{1}{W} & \gate{} \qwx[-1] & \multigate{1}{V} & \qw \\
& / \qw & \ghost{N} & \gate{} \qwx[-1] & \ghost{M} & \qw & \qw &  \gate{} \qwx[-1] & \qw & \qw & \ghost{W} & \gate{} \qwx[-1] & \ghost{V}  & \qw
}\nonumber
\end{eqnarray}

In the first step, according to Lemma~\ref{cosinesine}, for every $n$-qubit gate $U$, there exist $2^{n-1}$ 1-qubit gates $R_y(\theta_0)$, $R_y(\theta_1)$, $\cdots$, and $R_y(\theta_{2^{n-1}-1})$, four $(n-1)$-qubit gates $P_0$, $P_1$, $Q_0$, and $Q_1$, and three n-qubit gates $P = \outer{0}{0} \otimes P_{0} + \outer{1}{1} \otimes P_{1}$, $R = \sum^{2^{n-1}-1}_{i=0}R_y(\theta_i) \otimes \outer{i}{i}$, and $Q= \outer{0}{0} \otimes Q_{0} + \outer{1}{1} \otimes Q_{1}$, 
such that the first step is valid. 

In the second step, according to Lemma~\ref{ryrz}, for $R = \sum^{2^{n-1}-1}_{i=0}R_y(\theta_i) \otimes \outer{i}{i}$ gate, there exists one $n$-qubit gate $R_1 = \sum^{2^{n-1}-1}_{i=0}R_z(-\theta_i) \otimes \outer{i}{i}$, such that 
\begin{eqnarray*}
R & = & \sum^{2^{n-1}-1}_{i=0}R_y(\theta_i) \otimes \outer{i}{i} 
% \\ & = & 
\;\;=\;\;
\sum^{2^{n-1}-1}_{i=0} ( S^\dag \cdot H \cdot R_z(-\theta_i) \cdot H \cdot S )\otimes \outer{i}{i} \\
& = & \sum^{2^{n-1}-1}_{i=0} ( S^\dag \cdot H )\otimes \outer{i}{i} \cdot  \sum^{2^{n-1}-1}_{i=0}   R_z(-\theta_i) \otimes \outer{i}{i}  \cdot \sum^{2^{n-1}-1}_{i=0} ( H \cdot S )\otimes \outer{i}{i} \\
& = & (S^\dag \otimes I^{n-1}) \cdot (H \otimes I^{n-1}) \cdot R_1 \cdot (H \otimes I^{n-1}) \cdot (S \otimes I^{n-1})
\end{eqnarray*}

In the third step, according to Lemma~\ref{demultiplexing}, for $P = \outer{0}{0} \otimes P_{0} + \outer{1}{1} \otimes P_{1} $ gate, there exist $2^{n-1}$ 1-qubit gates $R_z(\alpha_0)$, $R_z(\alpha_1)$, $\cdots$, and $R_z(\alpha_{2^{n-1}-1})$, two $(n-1)$-qubit gates $V$ and $W$, and one n-qubit gate $R_0 = \sum^{2^{n-1}-1}_{i=0}R_{z}(\alpha_i) \otimes \outer{i}{i}$, such that the right-most 
change in the third step is valid. 
Similarly, according to Lemma~\ref{demultiplexing}, for $Q = \outer{0}{0} \otimes Q_{0} + \outer{1}{1} \otimes Q_{1} $ gate, there exist $2^{n-1}$ 1-qubit gates $R_z(\beta_0)$, $R_z(\beta_1)$, $\cdots$, and $R_z(\beta_{2^{n-1}-1})$, two $(n-1)$-qubit gates $M$ and $N$, and one n-qubit gate $R_2 = \sum^{2^{n-1}-1}_{i=0}R_{z}(\beta_i) \otimes \outer{i}{i}$, such that left-most change in the third step is valid.

We have arrived at a circuit that uses 
3 multiplexed $(n+1)$-qubit $R_z(\theta)$ gates,
4 $n$-qubit gates,
2 $H$ gates, and
4 $S$ gates
(because the $S^\dag$ gate can be written $S^\dag = SSS$).

We can now apply Lemma~\ref{rzcount} to 
the 3 multiplexed $(n+1)$-qubit $R_z(\theta)$ gates, and
we can apply the induction hypothesis to 
4 $n$-qubit gates.
The result is an implementation of $U$ built from 
$\{\C{X}, H, S, R_z(\theta)\}$.

Let $\mathcal{U}(n)$ denote the number of Clifford+R$_z$ gates
needed to implement an $n$-qubit unitary $U$.
We will show that
$\mathcal{U}(n) = 14 \times 4^{n-1} - 9 \times 2^{n-1}$.

In the base case of $n=1$, we have from 
Lemma~\ref{lem:clifford-plus-rx-is-universal-for-1-qubit-gates}
that
$\mathcal{U}(1) = 5$, and we see that
$\mathcal{U}(1) = 14 \times 4^{1-1} - 9 \times 2^{1-1} = 5$.

In the induction step, suppose the property holds for $n$,
where $n \geq 1$.
Consider the case of an $(n+1)$-qubit unitary $U$.
From the above calculation, we have
\(
\mathcal{U}(n+1) &=& 
4 \times \mathcal{U}(n) + 3 \times (2^{n+1} - 2 + 2^n) + 6
\;\;=\;\;
4 \times \mathcal{U}(n) + 9 \times 2^n
\)
where we use $S^\dag = SSS$.

Now we use the above equation and the induction hypothesis:
\begin{eqnarray}
\mathcal{U}(n+1)
& = &     4 \times \mathcal{U}(n) + 9 \times 2^n \nonumber
\;\;=\;\; 4 \times (14 \times 4^{n-1} - 9 \times 2^{n-1}) + 9 \times 2^n
\;\;=\;\;       14 \times 4^n - 9 \times 2^n \nonumber
\end{eqnarray}
which completes the proof.
\end{proof}

\mbox{} \\ \hspace*{1cm}
   [We need exponentially many gates to implement an $n$-qubit unitary in the worst case] \\
Notice that an $n$-qubit unitary $U$ has 
$2^n \times 2^n = (2^n)^2 = (2^2)^n = 4^n$ entries.
This makes the worst-case $14 \times 4^{n-1} - 9 \times 2^{n-1}$ gates 
needed to implement $U$ seem reasonable.

\mbox{} \\ \hspace*{1cm}
   [We prove one of the known lemmas] \\
Our next goal is to prove Lemma~\ref{rzcount}.
As a key step, we prove Lemma~\ref{rzrz} that decomposes 
an $(n+1)$-qubit multiplexed $R_z(\alpha)$ gate
into 
two $n$-qubit multiplexed $R_z(\alpha)$ gates 
plus
two $\C{X}$ gates.

\begin{lemma} {\bf (M{\"{o}}tt{\"{o}}nen et al.\ 2004)}
For $n \geq 1$,
let $R$ be the $(n+1)$-qubit gate 
$\sum_{i=0}^{2^n-1} R_z(\alpha_i) \otimes \outer{i}{i}$.
We can pick angles $\beta_i$ and $\gamma_i$ (for $0 \leq i < 2^{n-1}$) and
define 
$R_0 = \sum^{2^{n-1}-1}_{i=0}R_{z}(\beta_i) \otimes \outer{i}{i}$ and 
$R_1 = \sum^{2^{n-1}-1}_{i=0}R_{z}(\gamma_i) \otimes \outer{i}{i}$ 
such that 
the following circuit equality is correct:
\[ 
\Qcircuit @C=0.5em @R=0.4em {
  & \qw & \gate{R} \qwx[1] & \qw & &&  & & \qw & \targ & \gate{R_0}      & \targ & \gate{R_1} & \qw\\
 &  \qw &   \gate{} \qwx[1]     &  \qw   & &=& &   &   \qw &     \ctrl{-1}     &  \qw               &        \ctrl{-1}       & \qw  & \qw \\
 & \qw / & \gate{}        & \qw &  && & & \qw / & \qw       & \gate{} \qwx[-2] & \qw  &\gate{} \qwx[-2]    &\qw\\
}
\]
\label{rzrz}
\end{lemma}

Now we prove Lemma~\ref{rzcount}.

\begin{proof}[Proof of Lemma~\ref{rzcount}]
Let $\mathcal{R}_{\C{X}}(n)$ denote the number of $\C{X}$ gates 
needed to implement an $n$-qubit multiplexed $R_z(\alpha)$ gate.
Similarly, let $\mathcal{R}_{z}(n)$ denote the number of $R_z(\alpha)$ gates 
needed to implement an $n$-qubit multiplexed $R_z(\alpha)$ gate.
We will show that
$\mathcal{R}_{\C{X}}(n) = 2^n - 2$ and
$\mathcal{R}_{z}(n)    = 2^{n-1}$.

In the base case of $n=1$, we have that 
$\mathcal{R}_{\C{X}}(1) = 2^1 - 2 = 0$ and
$\mathcal{R}_{z}(1)    = 2^{1-1} = 1$, 
which is correct.

In the induction step, suppose the property holds for $n$, 
where $n \geq 1$.
Consider the case of an $(n+1)$-qubit multiplexed $R_z(\alpha)$ gate.
From Lemma~\ref{rzrz}, we have
\(
\mathcal{R}_{\C{X}}(n+1) &=& 2 \times \mathcal{R}_{\C{X}}(n) + 2 \\
\mathcal{R}_{z}(n+1)    &=& 2 \times \mathcal{R}_{z}(n)
\)
Now we use the two above equations and the induction hypothesis:
\begin{eqnarray}
\mathcal{R}_{\C{X}}(n+1) 
& = &     2 \times \mathcal{R}_{\C{X}}(n) + 2 
\;\;=\;\; 2 \times (2^n - 2) + 2 
\;\;=\;\; 2^{n+1} - 2 \nonumber
\\ 
\mathcal{R}_{z}(n+1) 
& = & 2 \times \mathcal{R}_{z}(n) 
\;\;=\;\; 2 \times 2^{n-1}
\;\;=\;\; 2^n \nonumber
\end{eqnarray}
which completes the proof.
\end{proof}


\newpage

\paragraph{Transition to MP3:}
Now we can finish the proof by focusing on approximation.

\paragraph{Main Point 3:}
We will prove the full theorem by approximating the $R_z$ gates. \\
\mbox{} \\ \hspace*{1cm}
   [We list properties of approximation of unitaries] \\

\begin{lemma}
\label{lem:hs-distance-to-itself-is-zero}
If $U \sim V$, then $d(U,V) = 0$. 
\end{lemma}

\begin{proof}
We can write $U = e^{i\alpha} V$, where $e^{i\alpha}$ is the global phase.
Notice $Tr( V^\dagger V ) = Tr( I_N ) = N$.
We calculate:
\(
d(U,V) &=& \sqrt{ 1 - \frac{ | Tr( U^\dagger V ) |^2 }{ N^2 } } 
 \;\;=\;\; \sqrt{ 1 - \frac{ | Tr( e^{-i\alpha} V^\dag V ) |^2 }{ N^2 } }
 \;\;=\;\; \sqrt{ 1 - \frac{ | N |^2 }{ N^2 } } 
 \;\;=\;\; 0
\)
\end{proof}

\begin{lemma}
\label{lem:hs-distance-and-global-phase}
If $U_1 \sim U_2$ and $V_1 \sim V_2$, then $d(U_1,V_1) = d(U_2,V_2)$.
\end{lemma}

\begin{proof}
From $U_1 \sim U_2$, we have that 
there exists $\alpha$ such that $U_1 = e^{i\alpha} U_2$.
From $V_1 \sim V_2$, we have that 
there exists $\beta$  such that $V_1 = e^{i\beta} V_2$.
We calculate:
\(
d(U_1,V_1) 
&=&
d(e^{i\alpha} U_2, e^{i\beta} V_2)
\\
&=&
\sqrt{ 1 - \frac{ | Tr( (e^{i\alpha} U_2)^\dagger (e^{i\beta} V_2) ) |^2 }{ N^2 } }
\\
&=&
\sqrt{ 1 - \frac{ | e^{-i\alpha} e^{i\beta} Tr( U_2^\dagger V_2 ) |^2 }{ N^2 } }
\\
&=&
\sqrt{ 1 - \frac{ | Tr( U_2^\dagger V_2 ) |^2 }{ N^2 } }
\\
&=&
d(U_2,V_2)
\)
\end{proof}

\begin{lemma}
$d(U_1 U_2, V_1 V_2) = d(V_1^\dag U_1, V_2 U_2^\dag)$.
\label{lem:distance-equality-for-products}
\end{lemma}

\begin{proof}
Notice that
\begin{eqnarray}
Tr((U_1 U_2)^\dag\ (V_1 V_2)) 
\;=\;
Tr(U_2^\dag U_1^\dag V_1 V_2) 
\;=\;
Tr(U_1^\dag V_1 V_2 U_2^\dag)
\;=\;
Tr((V_1^\dag U_1)^\dag V_2 U_2^\dag)
\label{eq:trace-equality-between-products}
\end{eqnarray}
In the first step, we use $(UV)^\dag = V^\dag\ U^\dag$.
In the second step, we use $Tr(UV) = Tr(VU)$.
In the third step, we use $(UV)^\dag = V^\dag\ U^\dag$.

Now we calculate:
\(
d(U_1 U_2, V_1 V_2) &=&
\sqrt{ 1 - \frac{ | Tr( (U_1 U_2)^\dag\ (V_1 V_2) ) |^2 }{ N^2 } }
\\ &=&
\sqrt{ 1 - \frac{ | Tr( (V_1^\dag U_1)^\dag\ (V_2 U_2^\dag) ) |^2 }{ N^2 } }
\\ &=&
d(V_1^\dag U_1, V_2 U_2^\dag)
\)
In the second step, we use Equation~\ref{eq:trace-equality-between-products}.
\end{proof}

\begin{lemma}
$d(I, V_1 V_2) = d(V_1^\dag, V_2)$ and
$d(U_1 U_2, I) = d(U_1, U_2^\dag)$.
\label{lem:distance-to-identity}
\end{lemma}

\begin{proof}
Immediate from Lemma~\ref{lem:distance-equality-for-products}.
\end{proof}

\begin{lemma}
$d(U,V) = d(V^\dag, U^\dag)$.
\label{lem:distance-of-conjugated-unitaries}
\end{lemma}

\begin{proof}
Immediate from Lemma~\ref{lem:distance-equality-for-products}.
\end{proof}

For approximation of a product of gates, errors add at most linearly.

\begin{lemma} {\bf (Wang and Zhang, 1994)}
$d(U,V) \leq d(U,I) + d(I,V)$.
\label{lem:trace-inequality}
\end{lemma}

\begin{lemma}
$d(U_1 U_2, V_1 V_2) \;\leq\; d(U_1, V_1) \;+\; d(U_2, V_2)$.
\label{lem:hs-small-product-rule}
\end{lemma}

\begin{proof}
We calculate:
\(
d(U_1 U_2, V_1 V_2) &=&
d(V_1^\dag U_1, V_2 U_2^\dag) 
\\ &\leq&
d(V_1^\dag U_1, I) + d(I, V_2 U_2^\dag)
\\ &=&
d(V_1^\dag, U_1^\dag) + d(V_2^\dag, U_2^\dag)
\\ &=&
d(U_1,V_1) + d(U_2,V_2)
\)
In the first step, we use Lemma~\ref{lem:distance-equality-for-products}.
In the second step, we use Lemma~\ref{lem:trace-inequality}.
In the third step, we use Lemma~\ref{lem:distance-to-identity}.
In the fourth step, we use Lemma~\ref{lem:distance-of-conjugated-unitaries}.
\end{proof}

\begin{lemma}
$d(\ U_1 \ldots U_n,\ \ V_1 \ldots V_n\ ) \;\leq\; \sum_{i=1}^n d(U_i,V_i)$.
\label{lem:hs-big-product-rule}
\end{lemma}

\begin{proof}
We proceed by induction on $n$.

In the base case of $n=1$, we have $d(U_1,V_1) \leq d(U_1,V_1)$, as required.

In the induction step, suppose the lemma holds for $n$.
Consider the case of $n+1$.
\(
d(\ U_1 \ldots U_{n+1},\ \ V_1 \ldots V_{n+1}\ ) 
&\leq&
d(\ U_1 \ldots U_n,\ \ V_1 \ldots V_n\ ) + d(U_{n+1},V_{n+1}) 
\\ &\leq&
\sum_{i=1}^n d(U_i,V_i) + d(U_{n+1},V_{n+1})
\\ &=&
\sum_{i=1}^{n+1} d(U_i,V_i) 
\)
In the first step, we use Lemma~\ref{lem:hs-small-product-rule}.
\end{proof}


\newpage

\hspace*{1cm}
   [We show how to approximate $R_z$ gates] \\
Recall that our goal is to replace each $R_z(\theta)$ gate by 
a circuit over $\{H,T\}$. 
Lemma~\ref{lem:approximation-of-rz} below states precisely that we can always
use $\{H, T\}$ to approximate $R_z(\theta)$, in the sense of using
$\{H, T\}$ to get as close as we like to any $R_z(\theta)$ gate,
% Indeed, $\{H, T\}$ is universal for defining $R_z(\theta)$ gates,

\begin{lemma}
For any real number $\theta$ and for $\epsilon > 0$,
there exists a circuit $C$ built from $\{H, T\}$,
such that $d(R_z(\theta),C) < \epsilon$.
\label{lem:approximation-of-rz}
\end{lemma}

We will prove Lemma~\ref{lem:approximation-of-rz} in three steps:
\begin{enumerate}
\item We express $R_z(\theta)$, up to a global phase, in terms of
real powers of two new gates $G_1$ and $G_2$ built from $\{H,T\}$.
\item We replace those real powers by nearby integer powers.
\item We prove that those integer powers provide arbitrarily close
approximations of the original real powers.
\end{enumerate}
%
Now we turn to the first step of the proof.
We define an intermediate gate set between $R_z(\theta)$ and
$\{H, T\}$, namely $\{G_1,G_2\}$, where
\(
G_1 &=& e^{-i\pi/4}\ THTH \\
G_2 &=& e^{-i\pi/4}\ HTHT
\)
At first glance, the choice of $G_1$ and $G_2$ may seem mysterious. 
The following observations explain why these gates are useful.
The idea behind $G_1$ and $G_2$ is that 
they represent rotations by irrational angles.
Intuitively, this means that their integer powers can come arbitrarily close to
any rotation about their respective axes.
Lemma~\ref{lem:properties-of-g1-g2} makes some observations about $G_1,G_2$.

\begin{lemma}
\label{lem:properties-of-g1-g2}
We have 
$\det(G_1) = \det(G_2) = 1$ and
$Tr(G_1) = Tr(G_2) = 1 + \frac{1}{\sqrt{2}}$ and
$G_1G_2 \neq G_2G_1$.
\end{lemma}

\begin{proof}
The proof is by straightforward matrix calculation.
In particular, we get:
\(
G_1G_2 &=& 
\frac{1}{2}
\left(
\begin{array}{cc}
1-i & -i\sqrt{2} \\
-i\sqrt{2} & 1+i
\end{array}
\right)
\quad \quad
\quad \quad
G_2G_1 \;\;=\;\;
\frac{1}{2}
\left(
\begin{array}{cc}
1-i\sqrt{2} & -i \\
-i & 1+i\sqrt{2}
\end{array}
\right)
\)
so $G_1G_2 \neq G_2G_1$.
\end{proof}

We use $G_1,G_2$ because both have determinant 1, like $R_z(\theta)$, and
they don't commute, which is one of the key properties needed for 
the generalized Euler decomposition below.

\begin{lemma}[\bf Nielsen and Chuang 2000, Exercise 4.8]
\label{lem:from-unitary-to-exponentiated-hamiltonian}
For any unitary $2 \times 2$ matrix $G$ such that $\det(G) = 1$,
there exists a Hermitian matrix $A$ and a real number $\alpha$
such that $A^2 = I$ and $G = e^{i\alpha A}$.
\end{lemma}

Lemma~\ref{lem:from-unitary-to-exponentiated-hamiltonian} lets us represent
$G_1,G_2$ by Hermitian matrices $A_1,A_2$ and real numbers $\alpha_1,\alpha_2$ 
such that
\begin{eqnarray}
A_1^2=I,
&&
A_2^2=I,
\label{eq:h1-squared-equals-h2-squared-equals-i}
\\
G_1 = e^{i\alpha_1 A_1}
&&
G_2 = e^{i\alpha_2 A_2}
\label{eq:g1-equals-the-exponentiation-of-h1-and-same-for-g2}
\end{eqnarray}
For the remainder of the proof, we fix particular choices of $A_1,A_2$ and 
$\alpha_1,\alpha_2$.
We don't need to know exactly what they are, but we will rely on
two key observations that are expressed by
Lemma~\ref{lem:h1-and-h2-are-not-scalar-multiples-of-one-another} and
Lemma~\ref{lem:trace-of-g1-and-trace-of-g2-expressed-via-alphas}.

\begin{lemma}
\label{lem:h1-and-h2-are-not-scalar-multiples-of-one-another}
$A_1$ and $A_2$ are not scalar multiples of one another.
\end{lemma}

\begin{proof}
If $A_1$ and $A_2$ were scalar multiples of one another,
then $G_1$ and $G_2$ would commute,
but Lemma~\ref{lem:properties-of-g1-g2} says that they don't.
\end{proof}

\begin{lemma}
\label{lem:trace-of-g1-and-trace-of-g2-expressed-via-alphas}
We have 
$Tr(G_1) = 2 \cos(\alpha_1)$ and
$Tr(G_2) = 2 \cos(\alpha_2)$.
\end{lemma}

\begin{proof}
Since $A_i$ is Hermitian and $A_i^2 = I$, its eigenvalues are $\{+1,-1\}$.
Therefore, the eigenvalues of $G_i = e^{i\alpha_i A_i}$ are
$e^{i\alpha_i}$ and $e^{-i\alpha_i}$,
from which we get
$Tr(G_i) = 2 \cos(\alpha_i)$.
\end{proof}

Lemma~\ref{lem:generalized-euler-decomposition} below is 
a variant of Euler decomposition of a unitary matrix $U$
with $\det(U) = 1$, in this case for our two nonparallel axes $A_1,A_2$ 
(Lemma~\ref{lem:h1-and-h2-are-not-scalar-multiples-of-one-another}).

\begin{lemma}[{\bf Nielsen and Chuang 2000, Exercise 4.11}]
% {\bf Segercrantz (1966)}]
\label{lem:generalized-euler-decomposition}
For any unitary $2 \times 2$ matrix $G$ such that $\det(G) = 1$,
there exist real numbers $a,b,c$ such that
$U \sim e^{iaA_1}\ e^{ibA_2}\ e^{icA_1}$.
\end{lemma}

For a unitary $G = e^{i\alpha A}$, where 
$\alpha$ is a real number and
$A$ is a Hermitian matrix,
we define $G^a = e^{i a\alpha A}$.
Notice that this definition only applies to matrices that are written in 
the form $e^{i \alpha A}$.
We can combine this definition with
Equation~\ref{eq:g1-equals-the-exponentiation-of-h1-and-same-for-g2} and write:
\begin{eqnarray}
G_1^a = e^{ia\alpha_1A_1},
\qquad
G_2^b = e^{ib\alpha_2A_2}.
\label{eq:g-1-a-and-g-2-b}
\end{eqnarray}

We can define $R_z(\theta)$ 
in terms of $\{G_1,G_2\}$, as expressed by Lemma~\ref{lem:from-rz-to-g1-g2}.

\begin{lemma}
\label{lem:from-rz-to-g1-g2}
For any $\theta$, there exist real numbers $a, b, c$ such that
$R_z(\theta) \sim G_1^a\ G_2^b\ G_1^c$.
\end{lemma}

\begin{proof}
We have that $R_z(\theta)$ is a unitary matrix and $\det(R_z(\theta)) = 1$. 
From Lemma~\ref{lem:generalized-euler-decomposition}, there exist real numbers
$x,y,z$ such that
\[
R_z(\theta)
\sim
e^{ixA_1}e^{iyA_2}e^{izA_1}.
\]

Equation~\ref{eq:g-1-a-and-g-2-b} says:
\[
G_1^a = e^{ia\alpha_1A_1},
\qquad
G_2^b = e^{ib\alpha_2A_2}.
\]
Thus, choosing
\[
a=\frac{x}{\alpha_1},
\qquad
b=\frac{y}{\alpha_2},
\qquad
c=\frac{z}{\alpha_1},
\]
we obtain
\[
R_z(\theta) 
\sim 
G_1^a\ G_2^b\ G_1^c.
\]
\end{proof}

If only the real numbers $a,b,c$ were integers, we would be done:
$G_1^a\ G_2^b\ G_1^c$ would be represented by a circuit built from $\{H,T\}$, 
up to a global phase.

The second step is to find integers $m,n,p$ that we can use 
instead of the real numbers $a,b,c$ in
Lemma~\ref{lem:from-rz-to-g1-g2} and 
get us as close as we like to $R_z(\theta)$.
The key lemma for this approximation is
Lemma~\ref{lem:approximate-g1-a-g2-b}, which is a special case of 
Lemma~\ref{%
lem:approximate-g-to-the-m-when-is-a-rotation-by-an-irr-multiple-of-pi}.
Intuitively, 
Lemma~\ref{%
lem:approximate-g-to-the-m-when-is-a-rotation-by-an-irr-multiple-of-pi}
talks about approximation based on a 1-qubit unitary $G$ that is a rotation by 
an irrational multiple of $\pi$.

\begin{lemma}[{\bf Hardy and Wright 2008}]
If the real number $\lambda$ is irrational, then 
\[
\{\,m\lambda \bmod 2 : m\in\mathbb Z\,\}
\]
is dense in $[0,2)$.
\label{lem:for-irrational-lambda-we-have-m-lambda-is-dense-mod-2}
\end{lemma}

\begin{lemma}
\label{lem:approximate-g-to-the-m-when-is-a-rotation-by-an-irr-multiple-of-pi}
For any $G = e^{i\lambda \pi A}$, where $\lambda$ is irrational and
$A$ is a Hermitian matrix that satisfies $e^{i 2 \pi A} = I$, 
and for any $\epsilon > 0$ and real number $a$, 
there exists an integer $m$ such that 
$d(G^a, G^m) < \epsilon$.
\end{lemma}

\begin{proof}
% Notice that $G^a = e^{i a \lambda \pi A}$.
% 
The function that maps $t$ to $e^{i t \pi A}$ is continuous,
so for our given $\epsilon$, we can find $\delta > 0$ such that
for all real numbers $t$, if $|a\lambda - t| < \delta$, then 
\(
d(e^{i (a \lambda) \pi A},\ e^{i t \pi A}) 
&<&
\epsilon
\)
Since $\lambda$ is irrational, we have from
Lemma~\ref{lem:for-irrational-lambda-we-have-m-lambda-is-dense-mod-2} 
that 
\[
\{\,m\lambda \bmod 2 : m\in\mathbb Z\,\}
\]
is dense in $[0,2)$.
Therefore, for our chosen $\delta$, there exist integers $m, k$ such that
\(
|a\lambda - (m \lambda + 2k)| &<& \delta
\)
from which we get
\(
d(e^{i (a \lambda) \pi A},\ e^{i (m \lambda + 2k) \pi A})
&<&
\epsilon
\)
Notice that
\(
e^{i (a \lambda) \pi A} 
&=&
G^a
\\
e^{i (m \lambda + 2k) \pi A}
&=&
e^{i (m \lambda) \pi A}
\;\;=\;\;
G^m
\)
so we conclude 
$d(G^a, G^m) < \epsilon$.
\end{proof}

\begin{lemma}
\label{lem:a-specific-lambda-is-irrational}
The real number $\frac{1}{\pi}
\arccos\!\left(
\frac{1}{2}
+
\frac{1}{2\sqrt{2}}
\right)$
is irrational.
\end{lemma}

\begin{proof}
Let
\[
\lambda
=
\frac{1}{\pi}
\arccos\!\left(
\frac12+\frac{1}{2\sqrt2}
\right).
\]
Suppose, for contradiction, that $\lambda$ is rational.
Then $e^{i\pi\lambda}$ is a root of unity, so it is an algebraic integer.
Hence
\[
e^{i\pi\lambda}+e^{-i\pi\lambda}
=
2\cos(\pi\lambda)
\]
is also an algebraic integer.

By the definition of $\lambda$,
\[
2\cos(\pi\lambda)
=
1+\frac{1}{\sqrt{2}}.
\]
Let
\[
x=1+\frac{1}{\sqrt{2}}.
\]
Then
\[
x-1=\frac{1}{\sqrt{2}},
\]
so
\[
2(x-1)^2=1.
\]
Thus
\[
2x^2-4x+1=0.
\]
The polynomial $2x^2-4x+1$ is irreducible over $\mathbb{Q}$, and it is
not monic. Hence $x$ is not an algebraic integer.

This contradicts the conclusion that
\[
2\cos(\pi\lambda)
\]
is an algebraic integer. Therefore $\lambda$ is irrational.
\end{proof}

\newpage

\begin{lemma}
\label{lem:approximate-g1-a-g2-b}
For any $\epsilon > 0$ and real number $a$,
there exists an integer $m$ such that 
$d(G_1^a, G_1^m) < \epsilon$.
Similarly, for any $\epsilon > 0$ and real number $b$,
there exists an integer $n$ such that 
$d(G_2^b, G_2^n) < \epsilon$.
\end{lemma}

\begin{proof}%[Proof of Lemma~\ref{lem:approximate-g1-a-g2-b}]
From 
Lemma~\ref{lem:properties-of-g1-g2} and
Lemma~\ref{lem:trace-of-g1-and-trace-of-g2-expressed-via-alphas},
we have
\(
Tr(G_i) 
= 2 \cos(\alpha_i) 
= 1 + \frac{1}{\sqrt{2}} 
\)
Now let us define 
\[
\lambda
=
\frac{1}{\pi}
\arccos\!\left(
\frac{1}{2}
+
\frac{1}{2\sqrt{2}}
\right)
\]
Hence, 
\(
\cos(\alpha_i) = \frac{1}{2} + \frac{1}{2\sqrt{2}} = \cos(\lambda \pi)
\)
Since $0 < \alpha_i < \pi$, 
we have
$\alpha_i=\lambda\pi$.

From Lemma~\ref{lem:a-specific-lambda-is-irrational}, 
we have that $\lambda$ is irrational.
From Equation~\ref{eq:h1-squared-equals-h2-squared-equals-i}, 
we get
\(
e^{i 2 \pi A_1} = I 
&&
e^{i 2 \pi A_2} = I
\)
Now the result follows from 
Lemma~\ref{lem:approximate-g-to-the-m-when-is-a-rotation-by-an-irr-multiple-of-pi}.
\end{proof}

% \begin{lemma}
% \label{lem:approximate-g2-b}
% For any real number $b$ and every $\epsilon > 0$, 
% there exists an integer $n$ such that \\
% $d(G_2^b, G_2^n) < \epsilon$.
% \end{lemma}
% 
% \begin{proof}%[Proof of Lemma~\ref{lem:approximate-g2-b}]
% We can show that $G_2$ is a rotation by an irrational multiple of $\pi$.
% Then the result follows from 
% Lemma~\ref{lem:approximate-g-to-the-m-when-is-a-rotation-by-an-irr-multiple-of-pi}.
% \end{proof}

The third step is to use
Lemmas~\ref{lem:from-rz-to-g1-g2} and \ref{lem:approximate-g1-a-g2-b}, and
to prove Lemma~\ref{lem:approximation-of-rz}.

\begin{proof}[Proof of Lemma~\ref{lem:approximation-of-rz}]
For a given $\theta$, 
we use Lemma~\ref{lem:from-rz-to-g1-g2} to find real numbers $a,b,c$ such that
\(
R_z(\theta) &\sim& G_1^a\ G_2^b\ G_1^c
\)
For a given $\epsilon > 0$, we use 
Lemma~\ref{lem:approximate-g1-a-g2-b} 
to find integers $m,n,p$ such that:
\(
d(G_1^a, G_1^m) &<& \epsilon/3 \\
d(G_2^b, G_2^n) &<& \epsilon/3 \\
d(G_1^c, G_1^p) &<& \epsilon/3
\)
Now we can calculate as follows:
\(
d(R_z(\theta), G_1^m\ G_2^n\ G_1^p) 
&=& 
d(G_1^a\ G_2^b\ G_1^c,\ G_1^m\ G_2^n\ G_1^p) 
\\
&\leq&
d(G_1^a, G_1^m) + d(G_2^b, G_2^n) + d(G_1^c, G_1^p)
\\
&<&
\epsilon/3 + \epsilon/3 + \epsilon/3
\;\;=\;\; 
\epsilon
\)
In the first step, we use Lemma~\ref{lem:hs-distance-and-global-phase}.
In the second step, we use Lemma~\ref{lem:hs-big-product-rule}.

If only the integers $m,n,p$ were nonnegative, we would be done:
$G_1^m\ G_2^n\ G_1^p$ would be represented by a circuit built from $\{H,T\}$,
up to a global phase.
If $m < 0$, then notice that 
$G_1^m = (G_1^{-1})^{-m}$, where 
\(
G_1^{-1} &=& G_1^\dag 
\;\;\sim\;\; 
(THTH)^\dag 
\;\;=\;\;
H^\dag\ T^\dag\ H^\dag\ T^\dag
\;\;=\;\;
H\ T^7\ H\ T^7
\)
Thus, also for $m < 0$, we have that $G_1^m$ is a circuit built from $\{H,T\}$.
If $n < 0$ or $p < 0$, we can handle those cases similarly.
We conclude that we can implement $G_1^m\ G_2^n\ G_1^p$ as a circuit built 
from $\{H,T\}$, 
up to a global phase.
\end{proof}

\newpage

\mbox{} \\ \hspace*{1cm}
   [We prove the full version of the theorem] \\
Now we can prove Theorem~\ref{thm:clifford-plus-t-is-universal}.

\begin{proof}[Proof of Theorem~\ref{thm:clifford-plus-t-is-universal}]
For a unitary $U$ and $\epsilon > 0$,
Lemma~\ref{lem:clifford-plus-rx-is-universal} says that
there exists a circuit $C_1$ built from $\{\C{X}, H, S, R_z(\theta)\}$,
such that $U \sim C_1$.

Let $k$ be the number of $R_z$ gates in $C_1$, and
let $\theta_1, \ldots, \theta_k$ be the angles used in those gates.
We have two cases.

In the first case,
if $k=0$, then $C_1$ is built from $\{\C{X},H,S\}$, and since
$S=T^2$, it is also built from $\{\C{X},H,T\}$.

In the second case, if $k>0$, then
from Lemma~\ref{lem:approximation-of-rz} we get circuits $W_j$ 
built from $\{H,T\}$ such that
\begin{eqnarray}
d(R_z(\theta_j),W_j) &<& \epsilon/k
\label{eq:approximation-of-r-z}
\end{eqnarray}
Let $C_2$ be a version of $C_1$ in which we have replaced
each $R_z(\theta_j)$ by $W_j$ and we have replaced each $S$ by $T^2$.
We calculate:
\(
d(U, C_2) &=& 
d(C_1, C_2) \;\;\leq\;\;
\sum_{j=1}^k d(R_z(\theta_j),W_j) \;\;<\;\;
k \cdot (\epsilon/k) \;\;=\;\; 
\epsilon
\)
In the first step, we use
Lemma~\ref{lem:hs-distance-and-global-phase}.
In the second step, we use 
Lemma~\ref{lem:hs-distance-to-itself-is-zero} and
Lemma~\ref{lem:hs-big-product-rule}.
In the third step, we use Equation~(\ref{eq:approximation-of-r-z}).
\end{proof}

We have seen have the proof of Theorem~\ref{thm:clifford-plus-t-is-universal}
rests on that irrational rotations generate dense sets of rotations about 
their axes. 
This mean that two suitably chosen axes are enough to generate 
all one-qubit rotations.

Overall, the proof relied on two ideas. 
First, the generalized Euler decomposition lets us express 
any one-qubit rotation in terms of rotations about two fixed axes. 
Second, because those rotations are by irrational angles, 
integer powers can approximate arbitrary real powers. 
Together, these ideas show that 
the gate set $\{H,T\}$ can approximate every $R_z(\theta)$ gate.

\paragraph{Transition to Close:} 
So there you have it.

\paragraph{Review:}
The key property of a universal gate set is that 
it can approximate any unitary.
Our proof strategy is to first prove the theorem for an easy gate set and then
extend to $\{\C{X}, H, T\}$.

\paragraph{Strong finish:}
Quantum computing has probabilities built in and
the idea of a universal gate set has approximation built in.
Quantum algorithms have no requirement to be exact; 
rather, they can try to get close.

\paragraph{Call to action:}
Implement the procedure for compiling any unitary into the easy gate set.

\end{document}

